\chapter{Scattering Kernels, LSZ, and Cluster Decomposition}
\label{ch:volume-ii-scattering-kernels}
\ConventionsLorentzian

\begin{definition}[Status of time-ordered correlators in the scattering chapters]
\label{def:scattering-time-ordered-correlator-status}
The time-ordered distributions used in this part have a fixed status.  In an
exact scattering statement, \(G_N^{\mathrm c}\) denotes the connected
time-ordered distribution associated with local Hilbert-space fields
\(\widehat\phi(x)\), obtained from the Wightman distributions by the specified
ordering and boundary-value prescription.  The input needed for LSZ is then
spectral: the two-point function has an isolated one-particle pole of mass
\(m\) and residue \(Z_\phi>0\), and the \(N\)-point time-ordered distribution
has the corresponding external pole boundary values after smearing.  No path
integral is part of the definition of the scattering operator in this
case; \(S\) remains \(\Omega_{\mathrm{out}}^*\Omega_{\mathrm{in}}\) from
Haag--Ruelle theory.

In perturbative examples, Feynman graphs, Feynman parameters, and loop
integrals denote coefficientwise contributions to a regulated source
functional \(Z[J]\).  Its connected functional is the coefficientwise
logarithm of \(Z[J]\), with the Lorentzian or Euclidean factor fixed by the
chosen regulator convention; the full source-functional status declaration is
given in Chapter~\ref{ch:volume-ii-generating-functionals-1pi}, and the BPHZ
status declaration is given in
Chapter~\ref{ch:volume-ii-subdivergences-forest-formulas}.  These are formal
or regulated calculations of the time-ordered distributions whose pole
residues may then be inserted into the LSZ formula.  Such graph calculations
do not define \(S\) and do not assert the existence of an unregulated
continuum path-integral measure.

If a Euclidean construction is used instead, the time-ordered Lorentzian
distributions are obtained only after the stated Osterwalder--Schrader
reconstruction or analytic-continuation hypotheses have been verified.  Thus a
symbol such as \(G_N^{\mathrm c}\), \(Z_\phi\), or a Feynman-parameter integral
always carries one of these statuses: exact Hilbert-space distribution,
Euclidean reconstructed distribution, or formal/regulator-dependent
perturbative coefficient.
\end{definition}

\begin{definition}[Mostly-plus Mandelstam variables]
\label{def:mostly-plus-mandelstam-variables}
Throughout the scattering and analyticity chapters, physical \(2\to2\)
momenta are written as incoming positive-energy momenta \(p_1,p_2\) and
outgoing positive-energy momenta \(p_3,p_4\), with
\[
  p_i^2=-m_i^2,
  \qquad
  p_1+p_2=p_3+p_4
\]
in the mostly-plus metric.  The Mandelstam variables are
\[
  s=-(p_1+p_2)^2,\qquad
  t=-(p_1-p_3)^2,\qquad
  u=-(p_1-p_4)^2 .
\]
They obey
\[
  s+t+u=m_1^2+m_2^2+m_3^2+m_4^2 .
\]
For equal masses \(m_i=m\), in the center-of-mass frame,
\[
  s=E_{\rm cm}^2,\qquad
  t=-2|\vec p_*|^2(1-\cos\theta),\qquad
  u=-2|\vec p_*|^2(1+\cos\theta),
\]
so physical \(s\)-channel scattering has \(s\ge4m^2\) and
\[
  -(s-4m^2)\le t\le0,\qquad
  -(s-4m^2)\le u\le0 .
\]
Positive real \(t\) or \(u\) in later fixed-\(t\) and crossing discussions is
therefore a crossed-channel timelike invariant, not a physical
\(s\)-channel scattering angle.
\end{definition}


\section{Scattering Operator and Kernel Notation}

This chapter fixes the distributional kernel notation used in the analytic
chapters.  The nonperturbative construction of \(S\) by Haag--Ruelle wave
operators was given in Chapters~\ref{ch:haag-ruelle} and
\ref{ch:volume-iv-haag-ruelle-mathematical-scattering}; it is not being
redefined here.

Work in one massive Haag--Ruelle scattering sector satisfying the hypotheses of
Theorem~\ref{thm:algebraic-haag-ruelle-scattering}, with the scalar
point-field construction and Cook estimates reviewed in
Proposition~\ref{prop:haag-ruelle-cook-limits}.  Fix a stable scalar
species of mass \(m>0\).  Its positive mass shell is
\[
  \Sigma_m^+
  =
  \{(\omega_{\vec p},\vec p):\vec p\in\mathbb R^d\},
  \qquad
  \omega_{\vec p}=\sqrt{\vec p^{\,2}+m^2}.
\]
The corresponding one-particle Hilbert space is denoted by \(\Hilb_1\).  In
the scalar case it is represented as \(L^2(\Sigma_m^+,\dd\mu_m)\), where
\[
  \dd\mu_m(\vec p)
  =
  \frac{\dd^d\vec p}{(2\pi)^d2\omega_{\vec p}} .
\]
The symmetric Fock space over \(\Hilb_1\) is
\[
  \mathcal F_s(\Hilb_1)
  =
  \bigoplus_{n=0}^\infty \Hilb_1^{\odot n}.
\]
The Haag--Ruelle wave operators from the preceding theorem layers
are isometries
\[
  \Omega_{\mathrm{in/out}}:
  \mathcal F_s(\Hilb_1)\longrightarrow \Hilb .
\]
On a scattering sector for which the incoming and outgoing ranges coincide,
the scattering operator is
\[
  S=\Omega_{\mathrm{out}}^*\Omega_{\mathrm{in}}
  :
  \mathcal F_s(\Hilb_1)\longrightarrow\mathcal F_s(\Hilb_1).
\]
This is the object whose kernels will be studied throughout this part.

Generalized momentum states on \(\Sigma_m^+\) are normalized by
\[
  \langle\vec q\,|\vec p\rangle
  =
  (2\pi)^d2\omega_{\vec p}\,\delta^{(d)}(\vec q-\vec p).
\]
For outgoing momenta \(q_i=(\omega_{\vec q_i},\vec q_i)\) and incoming
momenta \(p_j=(\omega_{\vec p_j},\vec p_j)\), write
\[
  S(q_1,\ldots,q_m|p_1,\ldots,p_n)
  =
  {}_{\mathrm{out}}\langle
  \vec q_1,\ldots,\vec q_m
  |\vec p_1,\ldots,\vec p_n
  \rangle_{\mathrm{in}} .
\]
This notation denotes a distributional kernel. Its value on wave packets
\(f_i,g_j\in C_c^\infty(\mathbb R^d)\) is
\[
\begin{aligned}
&S[f_1,\ldots,f_m;g_1,\ldots,g_n]
\\
&=
  \int
  \prod_{i=1}^m\dd\mu_m(\vec q_i)\,f_i(\vec q_i)^*
  \prod_{j=1}^n\dd\mu_m(\vec p_j)\,g_j(\vec p_j)
  S(q_1,\ldots,q_m|p_1,\ldots,p_n).
\end{aligned}
\]

\begin{center}
\begin{tikzpicture}[>=Stealth,scale=0.98]
  \draw[->,line width=0.75pt] (-4.2,-2.0)--(-4.2,2.25)
    node[above] {$x^0$};
  \draw[orange!85!black,dashed,line width=0.9pt] (-2.35,-1.18)--(2.25,-1.18);
  \draw[green!50!black,dashed,line width=0.9pt] (-2.35,1.18)--(2.25,1.18);
  \node[orange!85!black,anchor=west] at (2.4,-1.18)
    {incoming Fock basis};
  \node[green!50!black,anchor=west] at (2.4,1.18)
    {outgoing Fock basis};
  \draw[blue!75!black,line width=1.2pt,fill=blue!5] (0,0) ellipse (1.15 and 0.74);
  \node[blue!75!black] at (0,0) {\(S\)};
  \foreach \x in {-0.78,0,0.78} {
    \draw[orange!85!black,line width=1.05pt,->]
      (\x,-1.82)--({0.48*\x},-0.62);
  }
  \foreach \x in {-0.78,0,0.78} {
    \draw[green!50!black,line width=1.05pt,->]
      ({0.48*\x},0.62)--(\x,1.82);
  }
\end{tikzpicture}
\end{center}

\section{Connected Kernels and Invariant Amplitudes}

Let
\[
  P_{\mathrm{out}}=\sum_{i=1}^m q_i,
  \qquad
  P_{\mathrm{in}}=\sum_{j=1}^n p_j .
\]
Translation invariance implies that the non-identity part of every connected
scattering kernel is supported on \(P_{\mathrm{out}}=P_{\mathrm{in}}\).  The
phase convention is fixed once and for all by distinguishing the connected
\(S\)-kernel from the connected \(T\)-kernel.  With
\[
  S=\mathbf 1+\ii T,
\]
write
\[
\begin{aligned}
&S^{\mathrm c}(q_1,\ldots,q_m|p_1,\ldots,p_n)
\\
&\qquad =
  I^{\mathrm c}(q_1,\ldots,q_m|p_1,\ldots,p_n)
  +
  \ii\,T^{\mathrm c}(q_1,\ldots,q_m|p_1,\ldots,p_n),
\end{aligned}
\]
where \(I^{\mathrm c}\) is zero except for one incoming and one outgoing
particle of the same species, in which case
\[
  I^{\mathrm c}(q|p)
  =
  (2\pi)^d2\omega_{\vec p}\delta^{(d)}(\vec q-\vec p).
\]
The invariant amplitude \(\mathcal M\) is the reduced connected \(T\)-kernel:
\[
\begin{aligned}
&T^{\mathrm c}(q_1,\ldots,q_m|p_1,\ldots,p_n)
\\
&\qquad =
  (2\pi)^D\delta^{(D)}(P_{\mathrm{out}}-P_{\mathrm{in}})
  \mathcal M(q_1,\ldots,q_m|p_1,\ldots,p_n).
\end{aligned}
\]
Equivalently, the transition part of the connected \(S\)-kernel is
\[
  S^{\mathrm c}-I^{\mathrm c}
  =
  \ii(2\pi)^D\delta^{(D)}(P_{\mathrm{out}}-P_{\mathrm{in}})\mathcal M .
\]
The connected \(T\)-matrix kernel is therefore
\[
  {}_{\mathrm{out}}\langle q|T|p\rangle_{\mathrm c}
  =
  (2\pi)^D\delta^{(D)}(P_{\mathrm{out}}-P_{\mathrm{in}})
  \mathcal M(q|p).
\]

The word connected has a recursive distributional meaning. If the full
process splits into independent subprocesses, the full kernel is the product
of the kernels of those subprocesses, with separate momentum conservation in
each component. The connected kernel is the part that remains after subtracting
all such products.

\section{LSZ as Pole Extraction}

The pole-extraction theorem is Theorem~\ref{thm:lsz-wave-packet} of
Chapter~\ref{ch:lsz-reduction}; it is not restated here.  In this chapter we
use only its kernel consequence: after wave-packet smearing, a connected
time-ordered distribution whose external variables have isolated
one-particle poles of residue \(Z_\phi\) gives the connected scattering kernel
by multiplying each external leg by \(Z_\phi^{-1/2}\ii(k_a^2+m^2)\), taking
the simultaneous external pole boundary value, and then setting outgoing
variables \(k_i=q_i\) and incoming variables \(k_{m+j}=-p_j\) on
\(\Sigma_m^+\).  The resulting distribution is a matrix element of the
already-defined Haag--Ruelle operator
\(S=\Omega_{\rm out}^*\Omega_{\rm in}\); in perturbative sections the same
operation is applied coefficientwise to regulated time-ordered kernels.

\section{Cluster Decomposition}

Let \(\alpha=\{p_1,\ldots,p_n\}\) be the incoming labels and
\(\beta=\{q_1,\ldots,q_m\}\) the outgoing labels. A process partition
\(\Pi\) is a finite set of nonempty blocks \(B\), where each block contains a
subset \(\alpha_B\subset\alpha\) and a subset \(\beta_B\subset\beta\), and
the blocks partition all incoming and outgoing labels. Define
\[
  S_B^{\mathrm c}
  =
  S^{\mathrm c}(\beta_B|\alpha_B).
\]

\begin{proposition}[Connected scattering kernels]
\label{prop:connected-scattering-kernels}
For fixed incoming and outgoing labels, the full scattering kernel has a unique
expansion
\begin{equation}
  S(\beta|\alpha)
  =
  \sum_{\Pi}
  \prod_{B\in\Pi}S^{\mathrm c}(\beta_B|\alpha_B),
  \label{eq:scattering-kernel-cluster-partition}
\end{equation}
where the sum is over process partitions.  Equivalently, the connected kernel
is determined recursively by
\begin{equation}
  S^{\mathrm c}(\beta|\alpha)
  =
  S(\beta|\alpha)
  -
  \sum_{\Pi\ne\{\alpha\cup\beta\}}
  \prod_{B\in\Pi}S^{\mathrm c}(\beta_B|\alpha_B).
  \label{eq:connected-scattering-recursion}
\end{equation}
The one-particle identity kernel \(I^{\mathrm c}\) introduced above is included
as the connected kernel for a block with one incoming and one outgoing label:
\begin{equation}
  \begin{aligned}
  S^{\mathrm c}(q|p)
  &=
  I^{\mathrm c}(q|p)
  +
  \ii(2\pi)^D\delta^{(D)}(q-p)\mathcal M(q|p)
  \\
  &=
  (2\pi)^d2\omega_{\vec p}\delta^{(d)}(\vec q-\vec p)
  +
  \ii(2\pi)^D\delta^{(D)}(q-p)\mathcal M(q|p).
  \end{aligned}
  \label{eq:one-particle-connected-kernel}
\end{equation}
In a translationally invariant vacuum without external backgrounds, the
nontrivial one-particle amplitude vanishes for a stable isolated species, and
the identity term remains.
\end{proposition}

\begin{proof}
Order the finite label sets by their total number of external labels.  For one
block the statement defines \(S^{\mathrm c}\).  Suppose the connected kernels
have been defined for every proper subcollection.  Then
\eqref{eq:connected-scattering-recursion} defines the connected kernel for the
given collection by subtracting all factorizations into proper subprocesses.
Substituting this definition back into the right-hand side of
\eqref{eq:scattering-kernel-cluster-partition} gives \(S(\beta|\alpha)\);
uniqueness follows because every other solution would have a difference whose
highest-size nonzero component appears only in the one-block term and hence
must vanish.  This is the moment-cumulant recursion for scattering kernels,
with process partitions replacing ordinary set partitions.
\end{proof}

Proposition~\ref{prop:connected-scattering-kernels} is algebraic: it defines
connected kernels once the full kernel is known.  The physical cluster
property is the large-separation theorem below.  Its hypotheses are additional
Hilbert-space and locality hypotheses, not consequences of the recursion.

\begin{theorem}[Cluster decomposition for scattering kernels]
\label{thm:scattering-cluster-decomposition}
Assume a Haag--Ruelle scattering theory with a unique translation-invariant
vacuum, a mass gap above the vacuum, local commutativity, and the following
quantitative clustering estimate for the almost-local Haag--Ruelle
approximants used to construct the scattering states.  If \(F_T\) and \(G_T\)
are finite products of such approximants at fixed time parameter \(T\), and
if \(G_T(a)=U(a)G_TU(a)^*\), then for every \(N\) there are constants
\(C_N,M_N,R_N\), depending on the chosen operator families and wave packets
but independent of \(T\) and \(a\), such that
\begin{equation}
  \left|
  \bra{\vac}F_TG_T(a)\ket{\vac}
  -
  \bra{\vac}F_T\ket{\vac}\bra{\vac}G_T\ket{\vac}
  \right|
  \le
  C_N(1+|T|)^{M_N}(1+\rho(a,T))^{-N}.
  \label{eq:hr-quantitative-cluster-estimate}
\end{equation}
Here \(\rho(a,T)\) is the spacelike separation of the two localization tubes
after replacing the almost-local operators by local approximants; for
velocity supports separated from the light cone of the relative translation,
\(\rho(a,T)\ge c|a|-R_N(1+|T|)\) for some \(c>0\).  Let
\(\Psi_A^{\mathrm{in/out}}\) and \(\Psi_B^{\mathrm{in/out}}\) be two groups of
wave-packet scattering states whose velocity supports stay disjoint after a
relative spacelike translation \(a\), with \(a^2>0\) and \(|a|\to\infty\).
Then the full scattering matrix element factorizes:
\[
  \left\langle
    \Psi_A^{\mathrm{out}}\otimes U(a)\Psi_B^{\mathrm{out}},
    S\,
    \Psi_A^{\mathrm{in}}\otimes U(a)\Psi_B^{\mathrm{in}}
  \right\rangle
  \longrightarrow
  \left\langle\Psi_A^{\mathrm{out}},S\Psi_A^{\mathrm{in}}\right\rangle
  \left\langle\Psi_B^{\mathrm{out}},S\Psi_B^{\mathrm{in}}\right\rangle .
\]
Equivalently, the connected kernel
\(S^{\mathrm c}(\beta|\alpha)\) has no distributional component carrying
separate momentum conservation for a proper subprocess; it carries only the
total conservation law of the connected process.
\end{theorem}

\begin{proof}
Choose Haag--Ruelle approximants
\[
  \Psi_{A,T}^{\mathrm{out}}=A_T^{\mathrm{out}}\vac,\qquad
  \Psi_{A,T}^{\mathrm{in}}=A_{-T}^{\mathrm{in}}\vac,\qquad
  \Psi_{B,T}^{\mathrm{out}}=B_T^{\mathrm{out}}\vac,\qquad
  \Psi_{B,T}^{\mathrm{in}}=B_{-T}^{\mathrm{in}}\vac
\]
for the four scattering states.  The Haag--Ruelle theorem gives norm
convergence as \(T\to+\infty\) for the outgoing approximants and as
\(T\to+\infty\) for the incoming approximants written at time \(-T\).
Translation by \(U(a)\) is
unitary, so these norm errors are independent of the relative translation
\(a\).  Hence, for every \(\varepsilon>0\), there exists \(T_\varepsilon\)
such that for \(T\ge T_\varepsilon\) the exact matrix element differs by at
most \(\varepsilon\), uniformly in \(a\), from the finite-time expression
\[
  M_T(a)
  =
  \bra{\vac}
    (A_T^{\mathrm{out}})^*
    (B_T^{\mathrm{out}}(a))^*
    A_{-T}^{\mathrm{in}}
    B_{-T}^{\mathrm{in}}(a)
  \ket{\vac}.
\]
Almost-locality permits replacement of each approximant by a local operator
localized in its Haag--Ruelle tube, with an error decreasing faster than any
power of the tube size.  For the fixed \(T\) chosen above, large spacelike
\(a\) makes the \(A\)-tube spacelike separated from the translated
\(B\)-tube.  Local commutativity then orders the two groups, up to the
almost-local tail estimate, into
\[
  M_T(a)
  =
  \bra{\vac}F_TG_T(a)\ket{\vac}
  +
  O_N((1+|a|)^{-N}),
  \qquad
  F_T=(A_T^{\mathrm{out}})^*A_{-T}^{\mathrm{in}},
  \quad
  G_T=(B_T^{\mathrm{out}})^*B_{-T}^{\mathrm{in}} .
\]
Applying \eqref{eq:hr-quantitative-cluster-estimate} gives
\[
  \lim_{|a|\to\infty}M_T(a)
  =
  \bra{\vac}F_T\ket{\vac}\bra{\vac}G_T\ket{\vac}.
\]
Finally take \(T\to\infty\) in the two factors.  The same norm convergence
gives
\[
  \bra{\vac}F_T\ket{\vac}
  \longrightarrow
  \left\langle\Psi_A^{\mathrm{out}},S\Psi_A^{\mathrm{in}}\right\rangle,
  \qquad
  \bra{\vac}G_T\ket{\vac}
  \longrightarrow
  \left\langle\Psi_B^{\mathrm{out}},S\Psi_B^{\mathrm{in}}\right\rangle .
\]
Since \(\varepsilon\) was arbitrary, the scattering matrix element
factorizes in the stated large-\(|a|\) limit.

The kernel statement follows by testing the distributional kernel against
wave packets whose velocity supports realize the two translated clusters.  If
the connected kernel contained a component with separate momentum conservation
inside a proper subprocess, its wave-packet pairing would survive the
large-separation limit as a product contribution.  The recursion in
Proposition~\ref{prop:connected-scattering-kernels} subtracts exactly those
proper-subprocess products.  Therefore only the total momentum-conservation
delta distribution remains in \(S^{\mathrm c}\).
\end{proof}

\begin{center}
\begin{tikzpicture}[>=Stealth,scale=0.95,
  blob/.style={ellipse,draw=black,thick,minimum width=1.45cm,
    minimum height=0.95cm,fill=orange!8}]
  \node[blob] (A) at (0,0.72) {\(S^{\mathrm c}_1\)};
  \node[blob] (B) at (0,-0.72) {\(S^{\mathrm c}_2\)};
  \draw[thick,->] (-2.2,1.2)--(A.west);
  \draw[thick,->] (-2.2,0.45)--(A.west);
  \draw[thick,->] (A.east)--(2.2,1.0);
  \draw[thick,->] (-2.2,-0.45)--(B.west);
  \draw[thick,->] (B.east)--(2.2,-0.35);
  \draw[thick,->] (B.east)--(2.2,-1.2);
  \node[left] at (-2.25,1.2) {$p_1$};
  \node[left] at (-2.25,0.45) {$p_2$};
  \node[left] at (-2.25,-0.45) {$p_3$};
  \node[right] at (2.25,1.0) {$q_1$};
  \node[right] at (2.25,-0.35) {$q_2$};
  \node[right] at (2.25,-1.2) {$q_3$};
  \node at (3.15,0) {$\Longrightarrow$};
  \node[ellipse,draw=black,thick,minimum width=1.65cm,
    minimum height=1.3cm,fill=blue!6] (C) at (4.55,0) {\(S\)};
  \foreach \y in {1.05,0.35,-0.35} {
    \draw[thick,->] (3.65,\y)--(C.west);
  }
  \foreach \y in {1.05,0,-1.05} {
    \draw[thick,->] (C.east)--(5.85,\y);
  }
\end{tikzpicture}
\end{center}

The same decomposition appears in connected Green functions. Since LSZ acts
on external pole residues, it maps the connected part of a time-ordered Green
function to the connected part of the scattering kernel with the same
external labels.

\section{Preparation for Analytic Structure}

The object carried forward is the connected invariant amplitude
\[
  \mathcal M(q_1,\ldots,q_m|p_1,\ldots,p_n)
\]
as a distributional boundary value on the physical mass shells. Bound states
will appear as poles of this object in an appropriate channel. Resonances will
appear as poles reached by analytic continuation through a continuum cut.
Renormalization will be formulated as the construction of finite Green
functions and finite amplitudes while preserving the local form of the
counterterms.

For sectors with massless long-range forces, confinement, or infraparticle
spectra, the asymptotic data require modified observables or modified state
spaces. The later chapters return to those domains after the massive
scattering and renormalization structures have been developed. In particular,
one must not identify the charged electron labels used in perturbative QED
hard kernels with exact on-shell LSZ external particles of massless QED.  The
ordinary LSZ formula requires an isolated charged one-particle pole; the QED
infrared chapters instead construct resolution-inclusive probabilities and
softly dressed charged asymptotic data.
